\section{Answering the Research Questions}
\label{sec:rq_answers}
The following sections address each research question posed at the beginning of the study, mapping the empirical results to the original inquiries regarding framework design, configuration impact, and practical implementation bottlenecks under the finalized system architecture.

\subsection{RQ1: Designing a Modular Research-Grade Benchmarking Framework}
The project successfully demonstrated that a modular research-grade ZK-rollup benchmarking framework can be designed by isolating key performance factors into separate Rust microservices. The architecture splits the core pipeline into independent Sequencer, Executor, and Submitter components, connected via well-defined RPC interfaces. This modularity allows researchers to inject a benchmark-mode mock prover to abstract away full cryptographic overhead, or seamlessly integrate a real RISC0 prover to observe actual cycle limits. The architecture ensures that complex pipeline stages, such as Data Availability mode selection and Sequencer scheduling policies, can be toggled without requiring fundamental changes to the core state transition engine.

\subsection{RQ2: Impact of Configurations on Performance and Cost}
The empirical results clearly delineate the critical throughput-latency-cost trade-offs. 
\textbf{Batching:} Larger batch sizes significantly reduce per-transaction L1 gas costs by amortizing the fixed batch settlement and proof verification overhead. However, this efficiency comes at the direct expense of finalization latency, as early transactions wait longer in the sequencer pool.
\textbf{Sequencing:} Implementing a strict Nonce-Aware Priority Scheduler prevented out-of-order execution rejections. While the regular fee-priority scheduler maintained baseline throughput, the TimeBoost heuristic actively improved latency fairness under burst loads (achieving a Jain's Fairness Index of 0.77).
\textbf{Data Availability:} Shifting from standard calldata to EIP-4844 blobs resulted in approximately a 70\% reduction in per-transaction data publication costs, confirming blobs as a highly effective mechanism for scaling L2 footprint efficiency.

\subsection{RQ3: Practical Implementation Bottlenecks and Design Limitations}
Building and evaluating the reproducible prototype exposed several practical bottlenecks distinct from theoretical models. The most prominent bottleneck was cryptographic proof generation: integrating the real RISC0 prover revealed a massive fixed setup overhead (~70-80s) and strict step-function cycle limits governed by segment boundaries ($2^{20}$ cycles), which increased proving time by roughly 140-150s per segment. Additionally, reliance on an external Hardhat L1 network highlighted the volatility of finality; rapid block submissions occasionally caused nonce desynchronization, necessitating robust retry mechanisms (implemented via the Submitter's Saga pattern). Finally, the experiments highlighted that adaptive batching triggers are difficult to tune effectively under steady-state Poisson workloads, often causing artificial delays compared to strict maximum limits.
